\paragraph{生産者物価水準目標（PPLT）}
\label{sec:chapter5_beta_shock_policies_pplt}
式
\eqref{form:chapter3_policies_monetary_home_pplt_i_H_eq}、
\eqref{form:chapter3_policies_monetary_home_pplt_gamma_H_eq}および
\eqref{form:chapter3_policies_monetary_home_pplt_chi_H_eq}
により記述されたとおり、PPLTの式は以下のようであった。
\begin{align}
\label{form:chapter5_beta_shock_policies_pplt_i_H_eq}
    i_t^H&=i_{ss}^H+\phi_{gap}^H(\log(p_t^H)-\log(\chi_t^H))+\phi_{level}^H\gamma_t^H \\
\label{form:chapter5_beta_shock_policies_pplt_gamma_H_eq}
    \gamma_t^H&=\gamma_{t-1}^H+(\log(p_{t-1}^H)-\log(\chi_{t-1}^H)) \\
\label{form:chapter5_beta_shock_policies_pplt_chi_H_eq}
    \log(\chi_t^H)-\log(\chi_{ss}^H)&=\rho_{\chi}^H\left(\log(\chi_{t-1}^H)-\log(\chi_{ss}^H)\right)+\varepsilon_t^{\chi,H}
\end{align}
ここで
\(\phi_{gap}^H, \phi_{level}^H\)は反応係数、
\(p_t^H\)は生産者物価水準、
\(\chi_t^H\)は目標パス、
\(\gamma_t^H\)は過去の乖離の累積、
\(\varepsilon_t^{\chi,H}\)は外生ショック項である。
\begin{enumerate}
    \item
    \label{itm:chapter5_beta_shock_policies_pplt_c_H_W}
        \(c_t^{H\to W}\)：
        総消費指数\(c_t^{H\to W}\)の図
        \ref{fig:chapter5_beta_shock_graphs_c_H_W}
        を見ると、PPLTは消費者物価水準目標（CPLT）や消費者物価インフレ目標（IT-CPI）に見られるような
        極端な大暴落を免れていることがわかる。
        100期までの前半において、落ち込みの幅は
        総消費水準目標（CLT）や名目総消費水準目標（NCLT）には及ばないものの、
        ある程度の水準を維持している。消費の低下は厚生を悪化させる方向に働くため、
        この前半における厚生の押し下げはCPLTやIT-CPIよりも小さく防がれている。
        一方で中盤以降のプラス圏での推移を見ると、厚生を改善させる消費の増加幅は比較的小規模にとどまっている。
        CPLTやIT-CPIが遅れて形成する巨大な山や、NCLTが長期にわたって維持する高い水準と比較すると、
        後半の消費増加による厚生の押し上げ効果は限定的となっている。
    \item
    \label{itm:chapter5_beta_shock_policies_pplt_y_H}
        \(y_t^H\)：
        生産\(y_t^H\)の図
        \ref{fig:chapter5_beta_shock_graphs_y_H}
        を見ると、PPLTの軌道は総消費指数\(c_t^{H\to W}\)と同様に極端な変動を免れていることがわかる。
        100期までの前半において、PPLTの落ち込みは
        総消費水準目標（CLT）や名目総消費水準目標（NCLT）よりは大きいが、
        消費者物価水準目標（CPLT）や消費者物価インフレ目標（IT-CPI）のような極端な落ち込みには至っていない。
        生産の低下は労働の減少を意味するため厚生にはプラスに働くが、
        PPLTにおける前半の厚生押し上げ効果はCPLTやIT-CPIほど大きくはない。
        一方で中盤以降のプラス圏での推移（オーバーシュート）を見ると、
        生産の増加は労働の非効用を通じて厚生を悪化させる要因となる。
        PPLTのオーバーシュートはCLTやNCLTよりも低く抑えられており、
        またCPLTやIT-CPIのような遅くて巨大な山を作ることもないため、
        後半の生産過剰に伴う厚生のマイナス効果は比較的小さくとどまっている。
    \item
    \label{itm:chapter5_beta_shock_policies_pplt_pi_H}
        \(\pi_t^H\)：
        グロス・インフレ率\(\pi_t^H\)の図
        \ref{fig:chapter5_beta_shock_graphs_pi_H}
        を見ると、PPLTの定常状態からの乖離は生産者物価インフレ目標（IT-PPI）と並んで
        全政策の中で最も小さく抑えられていることがわかる。
        消費者物価インフレ目標（IT-CPI）のような初期の極端な下落や、
        総消費水準目標（CLT）に見られるような大きな変動は全く生じておらず、
        全期間を通じてゼロ近傍で極めて平坦な推移を維持している。
        生産者物価を直接の目標水準として強力にアンカーしているためインフレ率のブレがほぼ完全に制御されており、
        この定常状態からの乖離の小ささがフローの厚生損失を最小限にとどめることに大きく貢献している。
    \item
    \label{itm:chapter5_beta_shock_policies_pplt_Delta_t_minus_1_H}
        \(\Delta_{t-1}^H\)：
        価格分散\(\Delta_{t-1}^H\)の図
        \ref{fig:chapter5_beta_shock_welfare_Delta_H}
        を見ると、PPLTの定常状態からの乖離は
        生産者物価インフレ目標（IT-PPI）と並んで全政策の中で最も小さく抑えられていることがわかる。
        総消費水準目標（CLT）のような著しい増大や、
        消費者物価水準目標（CPLT）および消費者物価インフレ目標（IT-CPI）のような大きな山を作ることはなく、
        全期間を通じてほぼゼロの水準を維持している。
        これは前述のインフレ率の極めて安定した推移が価格分散の蓄積を完全に防いでいる結果であり、
        この生産効率の悪化の回避がPPLTの厚生を高く保つ重要な要因となっている。
\end{enumerate}
図\ref{fig:chapter5_beta_shock_graphs_p_H}が示すとおり、
PPLTは「水準目標」としての履歴依存性を持つため、
長期的には物価を定常状態へ戻す能力において全政策中で最も高い精度を見せる。
しかし、価格という「極めて動きの遅い敵」が動き出すのを待ってからしか反応できないという性質上、
ショック初期の数年間においては実体経済を深刻な不況にさらすこととなる。
厚生$W^H$の観点からは、長期的な物価の正確な復元よりも、初動でのジャンプによる需要の下支えが重要であり、
PPLTがCPLT同様に低い評価に留まる理由はここにある。